\section{Discussion and limitations}\label{sec:discussion}

\paragraph{What the guarantees cover.}
The linear-work results are constructive. They follow from a resource-incidence
model, an explicit temporal ordering and a certificate on the retained
hierarchy, and they bound the \emph{numerical} elimination for a fixed
regularized operator. Three costs are outside their scope. First, symbolic
setup: the current component-layout builder uses dense pair-position and
component-by-row tables of size $O(m_R^2+Jm_R)$ for $m_R=|R|$, so peak setup
memory is not yet linear (\cref{app:implementation}), and setup takes about as
long as the Newton loop at $N=100$ (\cref{tab:phases}). Second, the number of
outer iterations: we give no iteration bound for the implemented homogeneous
self-dual Mehrotra method, although the observed counts grow slowly. Third,
floating-point accuracy: exactness of condensation does not protect against
ill-conditioning, which is why every direction is checked against the full
operator and a fallback exists.

\paragraph{What the experiments show.}
The strongest results are the comparisons with MOSEK and BRISK
(\cref{sec:pushbox,sec:pendulum}), where SPRINT is faster from moderate horizons
on and, on PushBox, is the only solver that meets a strict $10^{-6}$ check on the
original SDP. They compare complete solver systems on different hardware (one
GPU versus a multicore CPU), and the baselines were measured earlier than the
current SPRINT version. The paired comparison with cuDSS (\cref{sec:cudss}) holds
the hardware and the IPM fixed. It shows a consistent but modest gain from the
hierarchy itself, which agrees with \cref{prop:ordering}: a general sparse
factorization can find comparable orderings, and the remaining advantage lies in
execution. The phase breakdown shows that the reduced system is a negligible
cost and that local partial factorization dominates, as \cref{thm:main}
predicts.

\paragraph{Future work.}
Four directions follow directly. (i) Assembling cone contributions directly into
local panels would remove the global sparse Schur matrix from the critical path.
(ii) Sparse indexing of the symbolic structures would extend the linear bounds
to total memory. (iii) Choosing the hierarchy by minimizing the cost
model~\eqref{eq:cost-model}, extended with hardware terms, would replace the
current empirical thresholds. (iv) Reusing factors across iterations, in the
spirit of inverse maintenance~\cite{jiang2022sos}, is attractive because spectral
approximation survives exact condensation without dependence on $N$
(\cref{prop:spectral}). However, a rank-one change of one scaling matrix can
induce a rank-$k$ change of the Schur matrix (\cref{prop:rank}), so cheap
updates are not automatic.

\section{Conclusion}\label{sec:conclusion}

Temporal repetition makes trajectory SOS relaxations large, but it also leaves
their coefficient resources and elimination interfaces local. We made this
precise with monomial separator bounds, an explicit linear-work ordering, and a
linear-time certificate under which interior elimination produces an exact
block-tridiagonal system with a border. The resulting cost bounds depend on
actual elimination fronts and interface widths, and they guided the design of a
GPU interior-point solver whose structured backend factors repeated local
patterns in batches and leaves only a small reduced system. In experiments, this
solver matches or exceeds the accuracy of MOSEK and BRISK and is faster from
moderate horizons on, by up to $3.2\times$ on PushBox and $12.3\times$ in solver
time on a 3D pendulum problem. We see this as evidence that interior-point
methods, combined with model-aware exact elimination on GPUs, can make
high-accuracy SOS relaxations practical for long-horizon robot trajectory
optimization. Because the analysis relies only on temporal locality, the same
approach applies to other multistage polynomial optimization problems, which we
plan to study next.
